% TOP_PROBLEM: {"id": 20001076, "title": "An order-statistic form of Shen's deficit conjecture", "category_id": 2, "category": {"id": 2, "name": "combinatorics", "display_name": "Combinatorics", "description": "Counting problems, graph theory, discrete structures.", "slug": "combinatorics", "order_index": 2, "created_at": "2026-07-31T15:26:25.671Z"}, "status": "partially_solved", "problem_number": "AIM-COMBINATORICS-0201", "set_id": 14, "statusQualification": "Uses simple digraphs and distinct-neighbor degrees, as required by this degree--girth conjecture; parallel-arc multiplicities would make the bound false."}
% FIRST_POSTED: 2026-10-01
% ENTRY_KIND: statement_only
% UnsolvedMath ID 20001076; source catalogue code AIM-COMBINATORICS-0201
% !TeX program = lualatex
% Definition and sources only; this file is not an LLM solution attempt.
\documentclass[11pt,a4paper]{article}
\usepackage[margin=25mm]{geometry}
\usepackage{fontspec}
\setmainfont{lmroman10-regular.otf}[BoldFont=lmroman10-bold.otf,
 ItalicFont=lmroman10-italic.otf,BoldItalicFont=lmroman10-bolditalic.otf]
\usepackage{amsmath,amssymb,mathtools}
\usepackage{unicode-math}
\setmathfont{latinmodern-math.otf}
\AtBeginDocument{\let\setminus\smallsetminus}
\usepackage{xurl,hyperref}
\hypersetup{colorlinks=true,urlcolor=blue}
\setcounter{secnumdepth}{0}
\setlength{\parindent}{0pt}
\setlength{\parskip}{5pt}
\setlength{\emergencystretch}{3em}
\begin{document}
\section*{An order-statistic form of Shen's deficit conjecture}\label{20001076}
{\small UnsolvedMath ID 20001076 · AIM-COMBINATORICS-0201 · Combinatorics · AIM Workshop Problem Lists}
\subsection{Definitions and mathematical statement}
Let $G$ be a finite simple loopless directed graph on $n$ vertices, with at most one arc from one vertex to another, and suppose every vertex has positive out-degree.
Write $d^+(v)=|\{w:v\to w\}|$.
Its directed girth $g$ is the least length of a directed cycle with distinct vertices; such a cycle exists because the finite graph has no vertex of out-degree zero.
Opposite arcs are allowed, in which case $g$ may equal two.

For any positive integer $r$, define the total deficit below out-degree $r$ by
\[
 t(G,r)=\sum_{v\in V(G)}\max\{0,r-d^+(v)\}.
\]
Vertices of degree at least $r$ contribute zero, so surplus degree does not cancel another vertex's deficit.
Shen's conjecture asks whether every such graph and every $r\geq1$ satisfy
\[
 n\geq r(g-1)+1-t(G,r).
\]
Equivalently, the total deficit should obey
$t(G,r)\geq r(g-1)+1-n$ whenever the right side is positive.
When the minimum out-degree is at least $r$, the deficit vanishes and the assertion recovers the familiar degree--girth inequality.
Here the question is whether the same obstruction persists with precisely one unit of allowed vertex-count loss for each missing outgoing neighbor.
No regularity condition is assumed, and there is no restriction on in-degrees.
Counting distinct outgoing neighbors, rather than arbitrary parallel arcs, is essential to the stated numerical bound.
\subsection{Short English statement}
Can the degree--girth bound be extended by subtracting exactly the total amount by which vertex out-degrees fall below a chosen threshold?
\subsection{Sources}
\begin{itemize}
\item[S1] ULAM.AI and the UnsolvedMath Contributors, supplied problems.json, record 20001076 (AIM-COMBINATORICS-0201), AIM Workshop Problem Lists. Catalogue identity and source transcription; CC BY 4.0.
\url{https://huggingface.co/datasets/ulamai/UnsolvedMath}.
\item[S2] American Institute of Mathematics, The Caccetta-Haggkvist conjecture, item 6. Original workshop question underlying AIM-COMBINATORICS-0201.
\url{https://aimath.org/WWN/caccetta/caccetta.pdf}.
\end{itemize}
\end{document}
